\documentclass[a4paper,12pt]{article}
\usepackage[utf8]{inputenc}
\usepackage[T1]{fontenc}
\usepackage{amsmath, amssymb, amsthm}
\usepackage{geometry}
\usepackage{xcolor}
\usepackage{graphicx}
\usepackage{fancyhdr}
\usepackage{enumitem}
\usepackage{ctex}
\usepackage[fontsize=14pt]{fontsize}
\usepackage{bm}



\setlist[itemize]{noitemsep, topsep=2pt}
\setlist[enumerate]{noitemsep, topsep=2pt}

\geometry{left=0.25cm, right=0.25cm, top=0.25cm, bottom=1.25cm}

% Define colors for highlighting
\definecolor{keywordcolor}{RGB}{0, 0, 200}
\definecolor{formulacolor}{RGB}{200, 0, 0}
\definecolor{examplecolor}{RGB}{0, 100, 0}

\newcommand{\key}[1]{\textcolor{keywordcolor}{\textbf{#1}}}
\newcommand{\formula}[1]{\textcolor{formulacolor}{\bfseries\boldmath{#1}}}
\newcommand{\formu}[1]{\textcolor{formulacolor}{\bm{#1}}}
\newcommand{\ex}[1]{\textcolor{examplecolor}{\bfseries\boldmath{#1}}}
\newcommand{\sol}[1]{\textcolor{examplecolor}{#1}}  
\renewcommand{\d}{\mathop{}\!\mathrm{d}}
\usepackage{multicol}

\usepackage{titlesec}

% 语法: \titlespacing*{<命令>}{<左间距>}{<上间距>}{<下间距>}
% 注意：<上间距> 和 <下间距> 可以使用 plus/minus 来设置弹性空间，
% 或者直接写固定数值（如 1em, 5pt）。

% 1. 调整一级标题 (\section)
\titlespacing*{\section}{0pt}{2pt}{2pt}

% 2. 调整二级标题 (\subsection)
\titlespacing*{\subsection}{0pt}{-2pt}{-2pt}

% 3. 调整三级标题 (\subsubsection)
\titlespacing*{\subsubsection}{0pt}{2pt}{2pt}

%分式使用行间版
\everymath{\displaystyle}


\title{\textbf{Differential Geometry: Key Definitions, Methods, and Examples}}
\author{Generated Summary}
\date{\today}

\begin{document}


\section{Curves in $\mathbb{R}^3$}

\subsection{Basic Definitions of curves}

\begin{itemize}
    \item \key{Length of a Curve:} The length of a parameterized curve $c$ is $\formu{L(c):=\int_a^b\|\dot{c}(t)\|\d t}$.
    \item \key{Regular Curve:} A curve $c(t)$ is regular if $\dot{c}(t) \neq 0$ for all $t\in I$.
    \item \key{Arc Length Parametrization}: A curve $c(s)$ is unit-speed if $\|\dot{c}(s)\| = 1$.\par
    For exerty regular curve, $\formu{\psi(t)=\int_{t_0}^t\|\dot{c}(t)\|\d t,\psi^{-1}\equiv \varphi,\tilde{c}=c\circ\varphi}$ is unit-speed.
    \item \key{Frenet Frame:} If $c$ is a unit-speed curve, $\formu{T(t):=\dot{c}(t),N(t):=\frac{\ddot{c}(t)}{\|\ddot{c}(t)\|}}$ has unit length nad are perpendicular to each other.
\end{itemize}

\begin{itemize}
    
    \item \key{Curvature ($\kappa$)}: Measures how fast the curve turns. For a unit-speed curve, $\formu{\kappa(t) = \|\ddot{c}(t)\|}$.\par
    Generally, $\formu{\kappa(t) = \frac{\|\dot{c}(t) \times \ddot{c}(t)\|}{\|\dot{c}(t)\|^3}=\frac{\sqrt{\|\ddot{c}(t)\|^2-\dot{v}(t)^2}}{v(t)^2}=\frac{\sqrt{v(t)^2\|\ddot{c}(t)\|^2-\left(\dot{c}(t)\cdot\ddot{c}(t)\right)^2}}{v(t)^3}}$
    \item \key{Torsion ($\tau$)}: Measures how much the curve twists out of the plane of curvature.\par
    For a general curve, $\formu{\tau(t) = \frac{(\dot{c}(t) \times \ddot{c}(t)) \cdot \dddot{c}(t)}{\|\dot{c}(t) \times \ddot{c}(t)\|^2}}$.
    \item \key{Frenet-Serret Frame}: Let $c$ be a unit-speed curve in $\mathbb{R}^3$, \formula{$T(t) = \dot{c}(t),N(t) = \frac{\ddot{c}(t)}{\left|\ddot{c}(t)\right|}=\frac{\ddot{c}(t)}{\left|\kappa(t)\right|},B(t)= T(t) \times N(t)$} are the tangent field, the normal field and the binormal field.$\{T, N, B\}$ is called the Frenet-Serret frame.
    \item \key{Frenet-Serret Formulas}:\formula{${\begin{pmatrix} T' \\ N' \\ B' \end{pmatrix} = \begin{pmatrix} 0 & \kappa & 0 \\ -\kappa & 0 & \tau \\ 0 & -\tau & 0 \end{pmatrix} \begin{pmatrix} T \\ N \\ B \end{pmatrix}}$}
\end{itemize}

\subsection{Extra Applications}

\begin{itemize}
    \item \key{Isoperimetric Inequality}: For a closed curve of length $L$ enclosing area $A$, \formula{${L^2 \geq 4\pi A}$} with equality iff the curve is a circle.
    \item \key{Winding Number}: Let $c$ a loop and $P_0$ be a point not lying on the curve, the winding number $W(c,P_0)$ is the number of times $c$ goes counter clockwise around $P_0$.\par
    If $c(t)=P_0+r(t)\left(\cos\theta(t),\sin\theta(t)\right)$ and $P_0=0$, then \formula{$W(c,P_0) = \frac{1}{2\pi}\int_a^b\frac{\dot{y}x-\dot{x}y}{x^2+y^2}(t) \, dt$}.
\end{itemize}



\subsection{Calculation Examples}
For a plane curve $c(t) = (x(t), y(t))$, the curvature is given by: \formula{$\kappa = \frac{|\dot{x}\ddot{y} - \dot{y}\ddot{x}|}{(\dot{x}^2 + \dot{y}^2)^{3/2}}$}.

\section{Surfaces in $\mathbb{R}^3$}

\subsection{Definition and Construction of Surfaces}
\begin{itemize}
    \item \key{Regular Map:} A map $\Phi$ is called regular if $\frac{\partial\Phi}{\partial x}(p)$ and $\frac{\partial\Phi}{\partial y}(p)$ are not zero and not parallel.
    \item \key{Coordinate Chart:} Given a subset $S\subseteq \mathbb{R}^3$, a coordinate chart on $S$ around a point $m\in S$, consists of
    \begin{itemize}
        \item an open subset $D\subseteq\mathbb{R}^2$
        \item a regular,injective smooth map $\Phi:D\to S$, $\Phi(D)\subseteq S$
        \item an open subset $\theta\in\mathbb{R}^3,m\in\theta$
    \end{itemize}
    such that $\Phi(D)=S\cap \theta$ and $\Phi^{-1}:\theta\cap S\to D$ is continuous.
    \item \key{Surface:} A subset $S\subseteq \mathbb{R}^3$ is a surface if every point $m\in S$ has a coordinate chart centered at $m$.\par
    There are three common surfaces:
    \begin{itemize}
        \item \key{Ruled Surface:} A ruled surface is a surface swept only by staight line $l:I\to\mathbb{R}^3$ moving along some curve $c:I\to\mathbb{R}^3\Rightarrow \formu{\Phi(t,r)=c(t)+r l(t),r\in\mathbb{R}}$ 
        \item \key{Graphs:} Evert smooth function of $D\subseteq\mathbb{R}^2\to\mathbb{R}^3$ determines a parameterized surface \formula{$S=\mathrm{graph}\;f=\left\{(x,y,z)\in D\times\mathbb{R}:z=f(x,y)\right\}$}
        \item \key{Level Sets:} Let $g:\mathbb{R}^3\to\mathbb{R}$ a smooth function.For every $c\in\mathbb{R}$, consider \formula{$S=g^{-1}(C)=\left\{(x,y,z)\in\mathbb{R}^3,g(x,y,z)=c\right\}$} is called a level set.\par
        If $\left(\frac{\partial g}{\partial x},\frac{\partial g}{\partial y},\frac{\partial g}{\partial z}\right) \neq 0$ on $g^{-1}(C)$, then $S$ is a surface.
    \end{itemize}
    \item \key{Smooth Map}: Given a surface $S$, a function $f:S\to\mathbb{R}^3$ is called smooth if all compositions $f\circ\Phi:\mathbb{R}^2\supseteq D\to\mathbb{R}^3$ are smooth for every coordinate chart $\Phi:D\to S$. 
\end{itemize}

\subsection{Vector Fields and Tangent Space}

\begin{itemize}
    \item \key{Tangent Vector:} A tangent vector $X$ at $m\in S$ is $X\in\mathbb{R}^3$ such that there exists a smooth curve $c:(-\epsilon,\epsilon)\to S$ with $c(0)=m$ and $\dot{c}(0)=X$.\par
    \item \key{Tangent Plane:} The set of all tangent vectors at $m\in S$ is called the tangent plane at $m$, denoted by $T_mS$.\par
    \ex{Cal Method:} Take the derivative of the surface $\nabla F$, and the tangent plane is \par
    \formula{$\left\{\omega\in\mathbb{R}^3:\nabla F\cdot w=0\right\}$} 
    \item \key{Vector Field:} A vecotr field on $S$ is a map $X:S\to\mathbb{R}^3$ such that $X(m)\in T_mS$ for all $m\in S$.
    \item \key{Directional Derivative}: Directional detivative of $f:S\to\mathbb{R}^k$ in the direction $X\in T_mS$ is defined by \formula{$T_mf(X)=\frac{\d}{\d t}\Bigl|_tf(m+tX)$}.
    \item \key{Vector field under local parametrization}: Given a coordinate chart $\Phi(u,v):D\to S$, a tangent vector $X$ on $S$ is the span of \formula{$\partial_u\Phi,\partial_v\Phi(\partial_u^\Phi,\partial_v^\Phi)$}.\par
    The tangent plane at $m=\Phi(u_0,v_0)$ is \formula{$T_m\Phi(x,y)=x\partial_u^\Phi(m)+y\partial_v^\Phi(m),(x,y)\in\mathbb{R}^2$}.
\end{itemize}

\subsection{Differential Forms}

\begin{itemize}
    \item \key{1-Form on Surface}: A 1-form on $S$ is a linear map $\alpha_m:T_mS\to \mathbb{R}$ for each point $m\in S$.We write $\alpha$ for the whole family. \formula{$\alpha_m(v)=\alpha_m(v_1\partial x+v_2\partial y)=v_1\alpha_m(\partial x)+v_2\alpha_m(\partial y)$}.\par
    1-form is called closed if $\d\alpha=0$ and is called exact if there exists a smooth function $f:S\to\mathbb{R}$ such that $\alpha=\d f$.exact implies closed.
    \begin{itemize}
        \item \key{1-Form $(\d x)_p,(\d y)_p$:} 1-form $(\d x)_p$ is defined by \formula{$(\d x)_p(\partial x)_p = 1,(\d x)_p(\partial y)_p=0$} and similarly for $(\d y)_p$. Moreover, every 1-form $\alpha$ at $p$ can be expressed as \formula{$\alpha=\alpha_1(p)(\d x)_p+\alpha_2(p)(\d y)_p$}, $\alpha_1,\alpha_2$ are smooth functions.
        \item \key{$T_pf$:} Given a smooth map $f:S\to\mathbb{R}$, the  directional derivative $T_pf$ is a 1-form at $p$ defined by \formula{$(\d f)_p(v)=T_pf(v)=(\nabla f)\cdot v$}.
    \end{itemize}
    \item \key{Local Representation of 1-Form}: Given a coordinate chart $\Phi(u,v):D\to S$, any 1-form $\alpha$ on $D$, then \formula{$\alpha^\Phi_{(x,y)}=\alpha_{\Phi(x,y)}\circ T_{(x,y)}\Phi$}.\par
    \ex{Cal Method:} If $\alpha=\alpha_1\d x+\alpha_2\d y$, then \ex{$\alpha^\Phi=\alpha_1(\Phi(u,v))\d x^\Phi+\alpha_2(\Phi(u,v))\d y^\Phi$}.
    \item \key{Integration of 1-Form along Curve}: Given a smooth curve $c:[a,b]\to S$ and a 1-form $\alpha$ on $S$, the integral of $\alpha$ along $c$ is defined by \formula{$\int_c\alpha:=\int_a^b\alpha_{c(t)}(\dot{c}(t))\d t$}.\par
    Speciallt, if $\alpha=\d f$ and $f$ is a smooth function, then \formula{$\int_c\d f=f(c(b))-f(c(a))$}.
    \item \key{2-Form:} A 2-form $\beta$ on $S$ is, for each point $m\in S$, a bilinear antisymmetric map $\beta_m:T_mS\times T_mS\to \mathbb{R}$.In particular, \formula{$\beta_m(v,v)=0$}.
    \item \key{k-Form:} A k-form $\omega$ on $S$ is, for each point $m\in S$, a multilinear antisymmetric map $\omega_m:(T_mS)^k\to \mathbb{R}$.
    \item \key{Wedge Product}: Given two 1-form $\alpha_1,\alpha_2$, their wedge product $\alpha_1\wedge\alpha_2$ is a 2-form defined by \formula{$(\alpha_1\wedge\alpha_2)_m(u,v)=(\alpha_1)_m(u)(\alpha_2)_m(v)-(\alpha_1)_m(v)(\alpha_2)_m(u),(u,v)\in T_mS$}.\par
    Note that $\alpha_1\wedge \alpha_2=-\alpha_2\wedge\alpha_1$.
    \item \key{Exterior Derivative of 1-Form}: $\alpha$ is a 1-form on $S$, its exterior derivative $\d\alpha$ is a 2-form defined by \formula{$(\d\alpha)_m\left(\partial^\Phi_x(m),\partial^\Phi_y(m)\right)=T_m\left(\alpha\left(\partial^\Phi_y\right)\right)\left(\partial^\Phi_x(m)\right)-T_m\left(\alpha\left(\partial^\Phi_x\right)\right)\left(\partial^\Phi_y(m)\right)$} for a chart $\Phi:D\to S$.\par
    \ex{Cal Method:} If $\alpha=\alpha_1\d u+\alpha_2\d v$, then \ex{$\d\alpha=\left(\frac{\partial \alpha_2}{\partial u}-\frac{\partial \alpha_1}{\partial v}\right)\d u\wedge \d v$}.
    \item \key{Properties of Exterior Derivative}:$f:S\to\mathbb{R}^3,\alpha$ 1-form, $\d(f\alpha)=\d f\wedge \alpha+f\d\alpha,\d(\d f)=0$
\end{itemize}

\subsection{Cartan's Moving Frame}

\begin{itemize}
    \item \key{Moving Frame}: A Moving frame on an open subset $U\subseteq\mathbb{R}^3$ consists of vector fields $e_1,e_2,e_3:U\to\mathbb{R}^3$ such that for every $x\in U$, $e_1(x),e_2(x),e_3(x)$ are an orthonormal basis of $\mathbb{R}^3$.\par
    \item \key{Connection 1-Form}: Given a moving frame $e_1,e_2,e_3$ on $S$, the connection 1-form $\omega_i^j$ is defined by \formula{$(\d e_i)_x(v)=\sum_{j=1}^3(\omega_i^j)_x(v)e_j(x)$}. In short, \formula{$\d e_i=\sum_{j=1}^3\omega_i^j e_j$}.For ONB, \formula{$(\omega_i^j)_x(v)=(de_i)_x(v)\cdot e_j(x)$}.
    \item \key{Connection Form and Frenet-Serret Frame:} For a unit-speed curve $c$ with $\kappa>0$.Suppose $e_1,e_2,e_3$ is an moving frame on $\mathbb{R}^3$ such that the restrictions of the fields to $c$ give the Frenet-Serret frame $T,N,B$.Then \formula{$(\omega_1^2)_{c(t)}\left(\dot{c}(t)\right)=\kappa,(\omega_1^3)_{c(t)}\left(\dot{c}(t)\right)=0,(\omega_2^3)_{c(t)}\left(\dot{c}(t)\right)=\tau$}.
    \item \key{Dual Frame:} Dual frames are 1-forms $\omega^k$ defined by \formula{$\omega^k(e_j)=\delta_j^k$}.\par
    \ex{Cal Method:} Given a chart $\Phi(u,v,w)$, then \ex{$\begin{cases}
        dx=\sum_{i=1}^3\omega^ie_i \\
        dx=\frac{\partial \Phi}{\partial u}\d u+\frac{\partial \Phi}{\partial v}\d v+\frac{\partial \Phi}{\partial w}\d w
    \end{cases}$}, so we can solve $\omega^i$.\par
    \item \key{First Cartan's Structure Equation}: \formula{$\d \omega^k + \sum_i \omega_i^k\wedge\omega^i=0$}.
    \item \key{Second Cartan's Structure Equation}: \formula{$\d \omega_i^j + \sum_k \omega_k^j\wedge\omega_i^k=0$}.
\end{itemize}

\subsection{Adapted Frame}
\begin{itemize}
    \item \key{Adapted Frame:} A moving frame $e_1,e_2,e_3:\mathcal{O}\to\mathbb{R}^3$ is called adapted to the surface $S$ if for every $p\in S\cap \mathcal{O}$, $e_3(p)$ is normal to $S$, i.e. $e_3(p)\times X=0$ for all $X\in T_pS$.\par
    In particular, if $S$ is oriented and $n$ is a unit normal field on $S$, then $n(p)=\pm e_3(p)$ for all $p\in S\cap\mathcal{O}$.
    $e_1(p),e_2(p)$ are a basis of $T_p(S),\forall\, p\in S\cap\mathcal{O}$.\par
    \ex{Cal Method:} If $\Phi(u,v)$ is the parametrization of $S$, \ex{$e_1=\frac{\Phi_u}{\|\Phi_u\|},e_3=\frac{\Phi_u\times\Phi_v}{\|\Phi_u\times\Phi_v\|},e_2=e_3\times e_1$}.\par
    For adapted frame, \formula{$\omega^3=0$ on $S$}.
    \item \key{Area Element under Adapted Frame}: If $e_1,e_2,e_3$ is an adapted frame on $S$, then the area element is given by \formula{$dA=\pm\omega^1\wedge\omega^2$}.\par
    \item \key{Positive:} An adapted frame is called positively oriented if $dA=\omega^1\wedge\omega^2\iff n=e_1\times e_2$.
\end{itemize}

\subsection{Guass Map and Shape Operator}
\begin{itemize}
    \item \key{Gauss Map:} If $n(p)$ has unit norm for $\forall\,p\in S$, then we call $n:S\to S^2\subseteq\mathbb{R}^3$ the Gauss map of $S$.Moreover, for the surface $\Phi:D\to S$, Gauss map $n$ can be writtrn as \formula{$n(u,v)=\frac{\Phi_u\times\Phi_v}{\|\Phi_u\times\Phi_v\|}$}.
    \item \key{Shape Operator (Weingarten Map)}: The shape operator at $p\in S$ is the linear map $\mathcal{S}_p:T_pS\to T_pS$ defined by \formula{$S_p(X)=-T_pn(X)$} for all $X\in T_pS$.\par
    If $e_1,e_2,e_3$ is an adapted frame on $S$, then \formula{$\mathcal{S}_p(v)=(\omega_1^3)(v)e_1(p)+(\omega_2^3)_p(v)e_2(p)$}.\par
    \ex{Cal Method:} If $\Phi(u,v)$ is the parametrization of $S$, then \ex{$\mathcal{S}_p(\Phi_u)= -n_u,\mathcal{S}_p(\Phi_v)=-n_v$}.
\end{itemize}


\subsection{First Fundamental Form and Second Fundamental Form}

\begin{itemize}
    \item \key{First Fundamental Form}: The first fundamental form at $p\in S$ is the bilinear map $I_p:T_pS\times T_pS\to \mathbb{R}$ defined by \formula{$I_p(v,w)=\langle v,w\rangle$} for all $v,w\in T_pS$.
    \item \key{Second Fundamental Form}: The second fundamental form at $p\in S$ is the bilinear map $II_p:T_pS\times T_pS\to \mathbb{R}$ defined by \formula{$II_p(v,w)=\langle v,\mathcal{S}_p(w)\rangle=I_p(v,\mathcal{S}_p(w))$} for all $v,w\in T_pS$.
    \item \key{Fundamental Forms under Local Parametrization}: Given a coordinate chart $\Phi(u,v):D\to S$, \formula{$E = \langle \Phi_u, \Phi_u \rangle, \quad F = \langle \Phi_u, \Phi_v \rangle, \quad G = \langle \Phi_v, \Phi_v \rangle,n = \frac{\Phi_u \times \Phi_v}{\|\Phi_u \times \Phi_v\|},e = \langle \Phi_{uu}, n \rangle, \quad f = \langle \Phi_{uv}, n \rangle, \quad g = \langle \Phi_{vv}, n \rangle$}.\par
    Then the first and second fundamental forms at $m$ are given by \formula{$I^\Phi=\begin{pmatrix} E & F \\ F & G \end{pmatrix},II^\Phi= \begin{pmatrix} e & f \\ f & g \end{pmatrix},\mathcal{S}=II\times I^{-1}=\frac{1}{EG-F^2}\begin{pmatrix} e & f \\ f & g \end{pmatrix}\begin{pmatrix} G & -F \\ -F & E \end{pmatrix}$}.
    \item \key{Area Element under Local Parametrization 2}: \formula{$dA = \sqrt{EG - F^2} \, du \, dv$}.
\end{itemize}



\subsection{Curvature}

\begin{itemize}
    \item \key{Normal Curvature:} The normal curvature in the $u-$direction is \formula{$\kappa_p(u)=II_p(u,u)=\langle u,\mathcal{S}_p(u)\rangle$}.
    \item \key{Principal Directions and Curvatures}: Principal curvature of $S$ at $p$ is the maximum and minimum values of $\kappa_p(u)$ for $\|u\|=1,u\in T_pS$.The corresponding directions are called principal directions.\par
    The principal curvatures and directions are the eigenvalues and eigenvectors of the shape operator $\mathcal{S}_p$.
    \item \key{Gaussian Curvature}: \formula{$K=k_1k_2$}.It is the determinant of the shape operator.
    \item \key{Mean Curvature}: \formula{$H=\frac{k_1+k_2}{2}$}.
    \item \key{Curvatures under Local Parametrization}: \formula{$K = \det(S) = \frac{eg - f^2}{EG - F^2},H = \frac{1}{2}\text{tr}(S) = \frac{eG - 2fF + gE}{2(EG - F^2)}$}
    The normal curvatures in direction $w = a \Phi_u + b \Phi_v$ is given by \formula{$\kappa\left(\frac{w}{\|w\|}\right) = \frac{e a^2 + 2f ab + g b^2}{E a^2 + 2F ab + G b^2}$}.Use $K$ and $H$ to find principal curvatures.
    \item \key{Umbilic Point}: A point $p\in S$ is called an umbilic point if \formula{$\kappa_1(p)=\kappa_2(p)$}\par
    Otherwise, \formula{$\kappa_1(p)k_2(p)>0$} is called elliptic point, \formula{$\kappa_1(p)k_2(p)<0$} is called hyperbolic point, and \formula{$\kappa_1(p)k_2(p)=0$} is called parabolic point.
    \item \key{Darboux Frame}: An adapted frame on $\mathbb{R}^3$ is called a Darbous frame if $e_1(p),e_2(p)$ are the principal directions of $S$ at $p$.
    \item \key{A Line of Curvature}: A unit-speed curve is called a line of curvature if $\dot{c}(t)$ is always a principal direction of $S$ at $c(t)$.
    \ex{Cal Method:} $\Phi$ is the parametrization of $S$, then a regular curve $c(t)=\Phi(x(t),y(t))$ is a line of curvature iff \ex{$\left(fE-eF\right)\dot{x}^2+\left(gE-eG\right)\dot{x}\dot{y}+\left(gF-fG\right)\dot{y}^2=0$}.
\end{itemize}

\subsection{Some equations of frames and curvatures}
For an adapted frame $e_1,e_2,e_3$ and its connnection forms and dual frames, we have
\begin{itemize}
    \item $\d\omega^1=\omega_1^2\wedge\omega^2,\d\omega^2=\omega_2^1\wedge\omega^1$
    \item $\omega_1^3\wedge\omega^1+\omega_2^3\wedge\omega^2=0$
    \item $\omega_1^3\wedge\omega^2+\omega^1\wedge\omega_2^3=2H\omega^1\wedge\omega^2$
    \item $\d\omega_1^2=-K\omega^1\wedge\omega^2,\d\omega_1^3=\omega_1^2\wedge\omega_2^3,\d\omega_2^3=\omega_2^1\wedge\omega_1^3$
    \item $\omega_1^3=II_{11}\omega^1+II_{12}\omega^2,\omega_2^3=II_{21}\omega^1+II_{22}\omega^2$
\end{itemize}\par
For Darboux frame, we have
\begin{itemize}
    \item $\d\kappa_1(e_2)=(\kappa_1-\kappa_2)\omega_1^2(e_1),\d\kappa_2(e_1)=(\kappa_1-\kappa_2)\omega_1^2(e_2)$
\end{itemize}

\subsection{Gauss Curvature and Theorema Egregium}
\begin{itemize}
    \item \key{Push Forward and Full back:} Let $F:\mathbb{R}^2\supseteq S\to\overline{S}\subseteq\mathbb{R}^3$
    \begin{itemize}
        \item \key{Push Forward of Curves:} For a curve $c:I\to S$, the push forward of $c$ by $F$ is \formula{$F_*c=F\circ c:I\to \overline{S}$}.
        \item \key{Push Forward of Vectors:} For a vector $v\in T_pS$, $c(t):I\to S$ s.t. $\dot{c}(0)=v$, the push forward of $X$ by $F$ is \formula{$F_*v=(F\circ c)(0)=T_PF(\dot{c}(0))$}, it's just the derivative $T_pF:T_pS\to T_{F(p)}\overline{S}$.
        \item \key{Pull Back of functions:} For a function $f:\overline{S}\to\mathbb{R}$, the pull back of $f$ by $F$ is \formula{$F^*f=f\circ F:S\to\mathbb{R}$}.
        \item \key{Pull Back of 1-Forms:} For a 1-form $\alpha$ on $\overline{S}$, the pull back of $\alpha$ by $F$ is \formula{$(F^*\alpha)_p(v)=\alpha_{F(p)}(T_pF(v))$} for all $p\in S,v\in T_pS$.
        \item \key{Pull Back of 2-Forms:} The pull back of 2-form $\beta$ by $F$ is \par
        \formula{$(F^*\beta)_p(u,v)=\beta_{F(p)}(T_pF(u),T_pF(v))$} for all $p\in S,u,v\in T_pS$.
        \item \key{Propositions:} If $\alpha,\beta$ are 1-forms on $\overline{S}$, then \formula{$F^*(\alpha\wedge\beta)=(F^*\alpha)\wedge(F^*\beta),F^*(\d\alpha)=\d(F^*\alpha)$}.
    \end{itemize}
    \item \key{Isometry}: A bijiective map $F:S\to\overline{S}$ is called an isometry if $F$ preserves the inner product of tangent vectors $\langle v,w\rangle=\langle T_pF(v),T_pF(w)\rangle,\forall\, v,w\in T_pS$, i.e. $F^*\bar{I}=I$. 
    \item \key{Local Isometry}: A map $F:S\to\overline{S}$ is called a local isometry if for every $p\in S$, there exists an open neighborhood $\mathcal{O}\subseteq S$ of $p$ such that the restriction $F|_{\mathcal{O}}:\mathcal{O}\to F(\mathcal{O})$ is an isometry.
    \item \key{Distance on Surface}: For $p,q\in S$, the distance between $p$ and $q$ on $S$ is defined by \formula{$d(p,q)=\inf\left\{L(c)\right\},c(0)=p,c(1)=q$}.\par
    \item \key{Isometry preserves distance}: If $F:S\to\overline{S}$ is an isometry, then \formula{$d(p,q)=\bar{d}(F(p),F(q))$} for all $p,q\in S$.
    \item \key{Theorema Egregium}: The Gauss curvature $K$ only depends on the first fundamental form $I$ on $S$ and it's invariant under local isometries $F:S\to\overline{S}$, i.e. \formula{$K(p)=\bar{K}(F(p))$} for all $p\in S$.
    \item \key{Behaviour of frame under isometry:} If $F:S\to\overline{S}$ is an isometry, and $e_1,e_2,e_3$ is an adapted frame on $S$, then \formula{$\bar{e}_i=F_*e_i=T_pF(e_i(p))$} is an adapted frame on $\overline{S}$. Dual frame $\{\bar{\omega}^1,\bar{\omega}^2\}$ is uniquely determined by \formula{$F^*\bar{\omega}^i=\omega^i$} for $i=1,2$. $\bar{\omega}^2_1$ is uniquely determined by \formula{$F^*\bar{\omega}^2_1=\omega^2_1$}.
    \item \key{Surface with all points umbilic}: Let $S$ be a connected surface such that all points are umbilic, $S$ has constant Gaussian curvature $K\geq 0$.And if $K=0$, then $S$ is a part of plane.If $K>0$, then $S$ is a part of sphere.
    \item \key{Constant Gaussian Curvature Surface:} Let $S$ be a compact surface in $\mathbb{R}^3$ with constant Gaussian curvature $K>0$, then $S$ is a sphere of radius $\frac{1}{\sqrt{K}}$. 
\end{itemize}

\subsection{Geodesics}

\begin{itemize}
    \item \key{Straight Line in Plane:} A curve $c(t)$ is a straight line $\iff \ddot{c}(t)=0 \iff \dot{c}(t)$ are parallel $\iff$ $c$ is the shortest path between two points.
    \item \key{Geodesic Curvature:} Let $\{e_1,e_2,e_3\}$ be an adapted frame such that $e_1(c(t))=\frac{\dot{c}(t)}{v}$, i.e. $e_1$ is the unit tangent vector of $c$, $e_3$ is the unit normal vector of $S$, $e_2=e_1\times e_3$,then the geodesic curvature of $c$ at $c(t)$ is defined by \formula{$\kappa_g(t)=\omega_1^2\left(\frac{\dot{c}(t)}{v}\right)=\omega_1^2(e_1)$}.\par
    \ex{Other Cal Method:} If $n$ is the unit normal vector of $S$, $r(t)$ is the parametrization of $c$, then \ex{$\kappa_g(t)=\frac{\det(\dot{r}(t),\ddot{r}(t),n(r(t)))}{\|\dot{r}(t)\|^3}=\frac{(\dot{r}(t)\times\ddot{r}(t))\cdot n(r(t))}{\|\dot{r}(t)\|^3}$}.\par
    \item \key{The relation between curvatures}: Let $\kappa(t)$ be the curvature of $c$ at $c(t)$ in $\mathbb{R}^3$, and $\kappa_n(t)$ be the normal curvature of $S$ at $c(t)$ in the direction of $\dot{c}(t)$, then \formula{$\kappa^2(t)=\kappa_g^2(t)+\kappa_n^2(t)$} and $\kappa_n(t)=\omega_1^3\left(\frac{\dot{c}(t)}{v}\right)=\omega_1^3(e_1)=\frac{c''\cdot n}{\|c'\|^2}$.
    \item \key{Geodesic}: A curve $c$ on $S$ is called a geodesic $\iff \kappa_g(t)=0 \iff \ddot{c}(t)$ is in the normal direction $\iff \omega_1^2(\dot{c}(t))=0\iff$ $c$ is locally the shortest path between two points on $S$.\par
    \ex{Cal Method:} If $\Phi(u,v)$ is the parametrization of $S$, then a regular curve $c(t)=\Phi(x(t),y(t))$ is a geodesic iff \ex{$\ddot{x}+\frac{1}{2E}(E_u\dot{x}^2+2E_v\dot{x}\dot{y}-G_u\dot{y}^2)=0$} and \ex{$\ddot{y}+\frac{1}{2G}(G_v\dot{x}^2+2G_u\dot{x}\dot{y}-E_v\dot{y}^2)=0$}.
    \item \key{Vector Field along Curve}: A vector field $X$ along a curve $c:I\to S$ is a map $X:I\to \mathbb{R}^3$ such that $X(t)\in T_{c(t)}S$ for all $t\in I$.
    \item \key{Covariant Derivative}: Let $e_1,e_2,e_3$ be a frame and $X(t)=f_1(t)e_1(t)+f_2(t)e_2(t)$ be a vector field along $c$. The covariant derivative of $X$ along $c$ is the vector field $\frac{\nabla}{\d t} X$ along $c$ defined by \formula{$\frac{\nabla}{\d t}X(t)=\left(\dot{f_1}+f_2\omega_2^1(\dot{c})\right)e_1\left(c(t)\right)+\left(\dot{f_2}+f_1\omega_1^2(\dot{c})\right)e_2\left(c(t)\right)$} for all $t\in I$.
    \item \key{Gauss-Weingarten Formula:} Let $e_1,e_2,e_3$ be an adapted frame on $S$, and $X(t)$ be a vector field along a curve $c:I\to S$. Then the derivative of $X$ in $\mathbb{R}^3$ can be decomposed as \formula{$\frac{\d}{\d t}X(t)=\frac{\nabla}{\d t}X(t)+II_{c(t)}\left(\dot{c}(t),X(t)\right)\cdot e_3\left(c(t)\right)$}.\par
    \item \key{Covariant Derivative and Geodesic:} Let $c$ be a unit-speed curve and $\{e_1,e_2,e_3\}$ an adapted frame with $e_1(c(t))=c(t)$, then \formula{$\frac{\nabla}{\d t}\left(\dot{c}(t)\right)=\kappa_g(t)e_2\left(c(t)\right),\kappa_n(t)=\kappa(\dot{c}(t))$}.\par
    Moreover, $c$ is a geodesic $\iff \frac{\nabla}{\d t}\left(\dot{c}(t)\right)=0$.
    \item \key{Parallel Transport:} Given a curve $c:I\to S$ and a vector $X_0\in T_{c(t_0)}S$, there exists a unique vector field $X$ along $c$ such that $X(0)=X_0$ and $\frac{\nabla}{\d t}X(t)=0$ for all $t\in I$.Such $X$ is called the parallel transport of $X_0$ along $c$.
    \item \key{Holomony Angle:} Given a curve $c:I\to S$ and the tangent frame $\{e_1,e_2\}$, the holomony angle of the parallel transport along $c$ is defined by \formula{$\mathrm{Hol}\,(c)=-\int_c\omega_1^2$}.
    \item \key{Energy of a Curve:} The energy of a curve $c:[a,b]\to S$ is \formula{$E(c)=\frac{1}{2}\int_a^b\|\dot{c}(t)\|^2\d t$}.
    \item \key{Exponential Map}: Given $p\in S$ and $v\in T_pS$, there exists a unique geodesic $c_v:I\to S$ such that $c_v(0)=p$ and $\dot{c}_v(0)=v$.The exponential map at $p$ is defined by \formula{$\exp_p:T_pS\to S,\exp_p(v)=c_v(1)$}.
\end{itemize}

\subsection{Fundamental Theorem of a Surface}
\begin{itemize}
    \item \key{Pullback of Fundamental Forms}: If $S$ is a parameterized surface via $\psi:\mathbb{R}^2\supseteq D\to S$, then $g=\begin{pmatrix}
        g_{11} & g_{12} \\ g_{12} & g_{22}
    \end{pmatrix}=\psi^*I,l=\begin{pmatrix}
        l_{11} & l_{12} \\ l_{12} & l_{22}
    \end{pmatrix}=\psi^*II$ where $g$ and $l$ is the families of symmetric metrics,parameterized by points $(x,y)\in D$.
    \item \key{Gauss-Codazzi Equations}: \formula{$\d\omega_1^2=\omega_1^3\wedge\omega_3^2,\d\omega_1^3=\omega_1^2\wedge\omega_2^3,\d\omega_2^3=\omega_2^1\wedge\omega_1^3$}.
    \item \key{Fundamental Theorem of a Surface}: Locally, a surface is uniquely determined by its first and second fundamental forms.
\end{itemize}

\subsection{Gauss-Bonnet Theorem}
\begin{itemize}
    \item \key{Gauss-Bonnet Formula for Compact Oriented serface:}\formula{$\frac{1}{2\pi}\int_S K\d A=2(1-\#\text{holes})=\chi(S)$}
    \item \key{Local version of Gauss-Bonnet Theorem:} Let $S$ be a compact oriented surface with the area element $\d A$.$\forall\, B\subseteq S$, whose boundary is smooth, admits an adapted, positive oriented frame, and $\partial B$ is a simple, unit-speed curve $\partial B:[a,b]\to S$, we have $\int_B K\d A+\int_a^b\kappa_g(t)\d t=2\pi$.
    \item \key{Euler Characteristic}: For a compact oriented surface $S$, the Euler characteristic is defined by \formula{$\chi(S)=V-E+F$}.$\chi(\text{Shpere})=2,\chi(\text{Torus})=0,\chi(\text{Disc})=1$
\end{itemize}

\section{Calculation Examples}

\subsection{Example: Sphere}
Parametrization: $\Phi(\theta,\varphi) = (R \sin \theta \cos \varphi, R \sin \theta \sin \varphi, R \cos \theta)$.\par
Adapted Frame: $e_1=$
\begin{itemize}
    \item \textbf{Coframe}: $\omega^1 = R d\theta$, $\omega^2 = R \sin\theta d\varphi$.
    \item \textbf{Differentials}:
    $d\omega^1 = 0$,
    $d\omega^2 = R \cos\theta d\theta \wedge d\varphi = \frac{\cos\theta}{R\sin\theta} \omega^1 \wedge \omega^2$.
    \item \textbf{Connection Form}:
    Assume $\omega_{12} = A \omega^1 + B \omega^2$.
    $0 = d\omega^1 = \omega^2 \wedge (-\omega_{12}) \implies \omega_{12} \propto \omega^2$.
    $d\omega^2 = \omega^1 \wedge \omega_{12}$.
    Result: \formula{$\omega_{12} = -\cos\theta d\varphi$}.
    \item \textbf{Curvature}: $d\omega_{12} = \sin\theta d\theta \wedge d\varphi = \frac{1}{R^2} \omega^1 \wedge \omega^2 \implies K = \frac{1}{R^2},H=-\frac{1}{R}$.
\end{itemize}

\subsubsection{Example: Cylinder}
Parametrization: $\Phi(\theta, z) = (r \cos \theta, r \sin \theta, z)$.
\begin{itemize}
    \item $E = r^2, F = 0, G = 1$.
    \item $e = -r, f = 0, g = 0$ (with inward normal).
    \item \textbf{Gaussian Curvature}: \formula{$K = \frac{(-r)(0) - 0}{r^2} = 0$}.
    \item \textbf{Mean Curvature}: $H = -1/(2r)$.
    \item \textbf{Geodesics}: Assume \sol{$\gamma(t)=x(t),y(t)$}, then \sol{$r^2\ddot{x}=0,\ddot{y}=0\Rightarrow x(t)=at+b,y(t)=ct+d$}.\par
\end{itemize}


\subsubsection{Example: Torus}
Parametrization: $\Phi(u, v) = ((R + r \cos v) \cos u, (R + r \cos v) \sin u, r \sin v)$.
\begin{itemize}
    \item \textbf{Gaussian Curvature}: \formula{$K = \frac{\cos v}{r(R + r \cos v)}$}
    \item Outer circle ($v=0$): $K > 0$. Inner circle ($v=\pi$): $K < 0$. Top/Bottom ($v=\pm \pi/2$): $K=0$.
\end{itemize}

\subsubsection{Example: Surface of Revolution}
Curve $c(t) = (r(t), 0, z(t))$ rotated around z-axis.
\begin{itemize}
    \item \textbf{Gaussian Curvature}: \formula{$K = -\frac{r''(t)}{r(t)}$} \quad \text{(assuming unit speed profile)}
\end{itemize}


\subsection{Algorithm: Computing Curvature via Coframe}

\begin{enumerate}
    \item \textbf{Choose a Frame}: Find orthonormal $e_1, e_2$ tangent to $S$.
    \item \textbf{Find Coframe}: Write $ds^2 = (\omega^1)^2 + (\omega^2)^2$. Identify $\omega^1, \omega^2$.
    \item \textbf{Differentiate}: Compute $d\omega^1$ and $d\omega^2$.
    \item \textbf{Solve for $\omega_{12}$}: Use the First Structure Equation:
    \begin{equation*}
        d\omega^1 = \omega^2 \wedge \omega_{21}, \quad d\omega^2 = \omega^1 \wedge \omega_{12}
    \end{equation*}
    (Recall $\omega_{21} = -\omega_{12}$).
    \item \textbf{Compute $K$}: Calculate $d\omega_{12}$ and write it as $-K \omega^1 \wedge \omega^2$.
\end{enumerate}


\subsection{The Darboux Frame (Frame along a Curve)}

For a unit-speed curve $c(s)$ on a surface $S$:
\begin{itemize}
    \item \key{Definition}: An orthonormal frame $\{T, V, n\}$ adapted to both the curve and the surface.
    \begin{itemize}
        \item $T(s) = \dot{c}(s)$ (Unit Tangent).
        \item $n(s)$ is the surface unit normal restricted to the curve.
        \item $S(s) = n(s) \times T(s)$ (Tangent Normal / Side Normal).
    \end{itemize}
    
    \item \key{Darboux Equations}\formula{$\begin{pmatrix} T' \\ S' \\ n' \end{pmatrix} = \begin{pmatrix} 0 & -\kappa_g & -\kappa_n \\ \kappa_g & 0 & -\tau_g \\ \kappa_n & \tau_g & 0 \end{pmatrix} \begin{pmatrix} T \\ S \\ n \end{pmatrix}$}
    
    \item \key{Curvatures}:
    \begin{itemize}
        \item \textbf{Geodesic Curvature} ($\kappa_g$): Bending in the tangent plane. $\kappa_g = \langle T', V \rangle = \omega_{12}(T)$.
        \item \textbf{Normal Curvature} ($\kappa_n$): Bending normal to the surface. $\kappa_n = \langle T', n \rangle = II(T, T) = \omega_{13}(T)$.
        \item \textbf{Geodesic Torsion} ($\tau_g$): Twisting of the normal. $\tau_g = \langle n', V \rangle = \omega_{23}(T)$.
    \end{itemize}
    
    \item \key{Relation to Space Curvature}:
    The curvature vector of the curve is $\vec{\kappa} = T' = \kappa_g V + \kappa_n n$.
    Squared curvature: \formula{$\kappa^2 = \kappa_g^2 + \kappa_n^2$}.
\end{itemize}


\section{Test Exam}
\subsection{EXERCISE 1. (True/False)}
\textit{Are the following claims true or false? Give a short one sentence explanation for your answer.}
\begin{enumerate}[label=(\roman*)]
    \item For any given smooth functions $\kappa>0$ and $\tau$ (defined on some interval) there exists a regular curve with curvature $\kappa$ and torsion $\tau$.\par
    \textbf{Answer: True.}
    \sol{Curvature and torsion uniquely determine a curve.}
    \item The subset of $\mathbb{R}^{3}$ defined by the equation $z=x^{2}-y^{2}$ is a surface.\par
    \textbf{Answer: True.}
    \sol{It can be seen as level set $g(x,y,z) = x^2 - y^2 - z=0$.$\nabla g = (2x, -2y, -1)\neq 0$.}

    \item Every surface admits a non-zero three-form.
    
    \textbf{Answer: False.}
    \sol{If $n>m$, then a n-form on an m-dimensional manifold is always zero. Since a surface is 2-dimensional, any 3-form on it must be zero.}

    \item The principal directions of a surface are intrinsic.
    
    \textbf{Answer: False.}
    \sol{The principal direction is the eigenvector of the shape operator, and the definition of the shape operator depends on the choice of the normal vector $n$ in the ambient space and the embedding, so the principal direction is an extrinsic geometric quantity.}

    \item If two isometric, connected surfaces in $\mathbb{R}^{3}$ have the same second fundamental form, then one is the image of the other under a composition of a translation, a rotation and a reflection.
    
    \textbf{Answer: True.}
    \sol{The Fundamental Theorem of Surfaces gets this.}

    \item If $v$ is a tangent vector at a point $p$ and $S_{p}$ the shape operator, then $S_{p}(v)$ is a vector normal to the surface.
    
    \textbf{Answer: False.}

    \item A curve $c$ in $\mathbb{R}^{3}$ of constant speed is a geodesic if and only if the acceleration is tangent to the surface $S$ at $c(t)$ for all $t$.
    
    \textbf{Answer: False.}
    \sol{On the contrary, the definition of a geodesic is that the tangent vector field is parallel transported along the curve, which means the acceleration vector $\ddot{c}(t)$ has no tangent component and must be \textbf{normal} to the surface (i.e., $\ddot{c}(t) \parallel n$).}
\end{enumerate}

\subsection{EXERCISE 2.}
\textit{Let $\theta, \varphi$ be the standard spherical coordinates on $S^{2}$, and consider the function $f:S^{2}\rightarrow\mathbb{R}$ given by $f(x,y,z)=(x-y)^{2}+z^{2}$. Calculate $\d f$ in terms of $\theta,\varphi$.}\par
The parameterization of a unit shpere is given by \sol{$\Phi(\theta,\varphi)=(\sin\theta\cos\varphi,\sin\theta\sin\varphi,\cos\theta)$}.\par
\textbf{Method 1:} Then we have \sol{$\setlength{\jot}{-4pt}\begin{aligned}[t]
    f &= (x-y)^2 + z^2 = x^2 - 2xy + y^2 + z^2 = (x^2 + y^2 + z^2) - 2xy \\
    &= 1 - 2(\sin\theta\cos\varphi)(\sin\theta\sin\varphi) \\
    &= 1 - 2\sin^2\theta \sin\varphi \cos\varphi = 1 - \sin^2\theta \sin(2\varphi)
\end{aligned}$}\par
Take the derivative, we get \sol{${\setlength{\jot}{-4pt}\begin{aligned}[t]
    df &= \frac{\partial f}{\partial \theta} d\theta + \frac{\partial f}{\partial \varphi} d\varphi \\
    &= \left( -2\sin\theta\cos\theta \sin(2\varphi) \right) d\theta + \left( -\sin^2\theta \cdot 2\cos(2\varphi) \right) d\varphi \\
    &= \formu{-\sin(2\theta)\sin(2\varphi) d\theta - 2\sin^2\theta\cos(2\varphi) d\varphi}
\end{aligned}}$}\par
\textbf{Method 2:} Alternatively, we can compute $\d f$ directly using the chain rule: \par
\sol{${\setlength{\jot}{-4pt}\begin{aligned}[t]
    df &= \frac{\partial f}{\partial x} dx + \frac{\partial f}{\partial y} dy + \frac{\partial f}{\partial z} dz = 2(x-y) dx + 2(y-x) dy + 2z dz \\
    &= 2(\sin\theta\cos\varphi - \sin\theta\sin\varphi)(\cos\theta\cos\varphi d\theta - \sin\theta\sin\varphi d\varphi) \\
    &\;\;\;\,+ 2(\sin\theta\sin\varphi - \sin\theta\cos\varphi)(\cos\theta\sin\varphi d\theta + \sin\theta\cos\varphi d\varphi) + 2(\cos\theta)(-\sin\theta d\theta) \\
    &=2\sin\theta(\cos\varphi-\sin\varphi)(\cos\theta(\cos\varphi-\sin\varphi)\d\theta-\sin\theta(\sin\varphi+\cos\varphi)\d\varphi)-2\cos\theta\sin\theta\d\theta \\
    &= \formu{-\sin(2\theta)\sin(2\varphi) d\theta - 2\sin^2\theta\cos(2\varphi) d\varphi}
\end{aligned}}$}

\subsection{EXERCISE 3.}
\textit{The catenoid $S$ is given by $\Psi(u,v)=(\cosh v \cos u, \cosh v \sin u, v),u\in [0,2\pi),v\in\mathbb{R}$.}
\begin{enumerate}[label=(\roman*)]
    \item \textbf{Determine an adapted moving frame $\{e_1,e_2,e_3\}$ on $S$.}\par
    Take partial derivatives: \sol{$\Psi_u = (-\cosh v \sin u, \cosh v \cos u, 0),\Psi_v = (\sinh v \cos u, \sinh v \sin u, 1)$}.\par
    These two vectors are orthogonal, thus we can unitize them to get the tangent frame $e_1,e_2$. \sol{$
    e_1 = \frac{\Psi_u}{\|\Psi_u\|} = \frac{\Psi_u}{\cosh v} = \formu{(-\sin u, \cos u, 0)},e_2 = \frac{\Psi_v}{\|\Psi_v\|} = \frac{\Psi_v}{\sqrt{\sinh^2 v + 1}} = \frac{\Psi_v}{\cosh v} = \formu{(\tanh v \cos u, \tanh v \sin u, \mathrm{sech} v)}$}\par
    $e_3=e_1\times e_2:$, \sol{$e_3 = \begin{vmatrix} i & j & k \\ -\sin u & \cos u & 0 \\ \tanh v \cos u & \tanh v \sin u & \mathrm{sech} v \end{vmatrix} 
    = \formu{(\mathrm{sech} v \cos u, \mathrm{sech} v \sin u, -\tanh v)}$}
    \item \textbf{Calculate the dual frame $\{\omega^1,\omega^2\}$ and the connection forms $\omega_i^j$.}\par
    Since $e_1=\frac{\partial_u}{\cosh v}$.And $\omega^1(e_1)=1$, thus \sol{$\omega^1 = \cosh v du$}.The same reason to the \sol{$\omega^2=\cosh v\d v$}.\par
    Another solution is $\begin{cases}
        \d x=\omega^1 e_1 + \omega^2 e_2 \\
        \d x=\Phi_u du + \Phi_v dv= \cosh v e_1 du + \cosh v e_2 dv
    \end{cases}$ and we can get the same result.\par
    $\omega_1^2=\d e_1\cdot e_2=\left(-\cos u\d u,-\sin u\d u\right)\cdot\left(\tanh v \cos u, \tanh v \sin u, \mathrm{sech} v\right)=-\tanh v\cos^2 u\d u-\tanh v\sin^2 u\d u=\formu{-\tanh v\d u}$.\par
    The same reason to get \formula{$\omega_1^3=-\mathrm{sech}\,v\d u,\omega_2^3=-\mathrm{sech}\,v\d v$}.
    \item \textbf{Calculate the mean curvature $H$ and the Gaussian curvature $K$ of $S$.}\par
    \ex{$d\omega_1^2 = -K \omega^1 \wedge \omega^2$} and $d\omega_1^2 = d(-\tanh v \, du) = -\mathrm{sech}^2 v \, dv \wedge du = \mathrm{sech}^2 v \, du \wedge dv,\omega^1 \wedge \omega^2 = \cosh^2 v \, du \wedge dv$.\par
    Campare the coefficients, we get \formula{$K = -\frac{1}{\cosh^4 v}$}.\par
    For mean curvature $H$, we have \ex{$\omega_{13} \wedge \omega^2 + \omega^1 \wedge \omega_{23} = 2H \omega^1 \wedge \omega^2$}.\par
    Substitute the values, we get \formula{$H=0$}.
    \item \textbf{Verify that the total Gaussian curvature $\int_S K dA$ of $S$ is $-4\pi\left(\frac{\d}{\d v}\tanh v=\frac{1}{\cosh^2v}\right)$.}\par
    \ex{$\d A=\omega^1\wedge \omega^2=\cosh^2 v \, du \wedge dv$}.\par
    Thus we have $=\begin{aligned}[t]
    \int_S K dA &= \int_{-\infty}^{\infty} \int_{0}^{2\pi} \left(-\frac{1}{\cosh^4 v}\right) (\cosh^2 v) \, du \, dv \\
    &= \int_{-\infty}^{\infty} \int_{0}^{2\pi} -\frac{1}{\cosh^2 v} \, du \, dv = 2\pi \int_{-\infty}^{\infty} -\mathrm{sech}^2 v \, dv \\
    &= -2\pi \left[ \tanh v \right]_{-\infty}^{\infty} = -2\pi (1 - (-1)) = \formu{-4\pi}
    \end{aligned}$
\end{enumerate}



\subsection{EXERCISE 4.}
\textit{Assume the first fundamental form od a surface $S$ in a certain parametrization is given by $I_{ij}=\begin{pmatrix}
    1/y^2 & 0 \\ 0 & 1/y^2
\end{pmatrix}$ on the half-plane $D=\{y>0\}$.}
\begin{enumerate}[label=(\roman*)]
    \item \textbf{Calculate a tangent frame field $e_1,e_2$ on $D$.}\par
    Using the first fundamental form, we get \sol{$\langle \partial_x,\partial_x\rangle= \frac{1}{y^2}$}, thus $\|\partial_x\|=\frac{1}{y}$, so we have \formula{$e_1=y\partial_x$}.\par
    Similarly, we get \formula{$e_2=y\partial_y$}.
    \item \textbf{Determine the connection form $\omega_1^2$ uniquely defined by the structure equation.}\par
    Firstly, we can get the dual frame \sol{$\omega^1 = \frac{1}{y} dx, \omega^2 = \frac{1}{y} dy$}.\par
    Assume $\omega_1^2=A\d u+B\d v$, then we have \ex{$d\omega^1 = \omega_1^2 \wedge \omega^2$} and \ex{$d\omega^2 = \omega^1 \wedge \omega_1^2$}.\par
    $\d\omega^1=-\frac{1}{y^2}\d y\wedge\d x=\frac{1}{y^2}\d x\wedge\d y=(A\d x+B\d y)\wedge\frac{1}{y}\d y=\frac{A}{y}\d x\wedge\d y\Rightarrow$\ex{$A=\frac{1}{y}$}.\par
    $\d\omega^2=-\frac{1}{y^2}\d y\wedge\d y=0=(\frac{1}{y}\d x)\wedge(A\d x+B\d y)=\frac{B}{y}\d x\wedge\d y\Rightarrow$\ex{$B=0$}.\par
    Thus we have \formula{$\omega_1^2 = \frac{1}{y} dx$}.
    \item \textbf{Determine the geodesics.}\par
    The first fundamental form is \sol{$I_{ij} = \begin{pmatrix} E & F \\ F & G \end{pmatrix} = \begin{pmatrix} 1/y^2 & 0 \\ 0 & 1/y^2 \end{pmatrix}$}.We know that a regular curve $\gamma(t) = (x(t), y(t))$ on the surface is a geodesic if it satisfies the following system of differential equations: \ex{$\ddot{x} + \frac{1}{2E}(E_x\dot{x}^2 + 2E_y\dot{x}\dot{y} - G_x\dot{y}^2) = 0,\ddot{y} + \frac{1}{2G}(G_y\dot{x}^2 + 2G_x\dot{x}\dot{y} - E_y\dot{y}^2) = 0$}.\par
    Given $E = y^{-2}$, $G = y^{-2}$, and $F=0$, the partial derivatives are: $E_x = \frac{\partial}{\partial x}(y^{-2}) = 0,E_y = \frac{\partial}{\partial y}(y^{-2}) = -2y^{-3},G_x = \frac{\partial}{\partial x}(y^{-2}) = 0,G_y = \frac{\partial}{\partial y}(y^{-2}) = -2y^{-3}$.\par
    Substituting into the first equation, we have \sol{$\ddot{x}+\frac{y^2}{2}\left( 0 \cdot \dot{x}^2 + 2(-2y^{-3})\dot{x}\dot{y} - 0 \cdot \dot{y}^2 \right)=\ddot{x}-\frac{2\dot{x}\dot{y}}{y}=0$}.\par
    And the same to the second equation, \sol{$\ddot{y}+\frac{y^2}{2}\left( (-2y^{-3})\dot{x}^2 + 2(0)\dot{x}\dot{y} - (-2y^{-3})\dot{y}^2 \right)=\ddot{y}-\frac{1}{y}(\dot{x}^2 - \dot{y}^2)=0$}.\par
    \textbf{Case 1:} $\dot{x}=0$, then $x(t)=x_0,\ddot{x}=0$, therefore $\ddot{y}+\frac{1}{y}\dot{y}^2=0\Rightarrow y=\sqrt{At+B}$.\par
    \textbf{Case 2:} $\dot{x}\neq 0$, from the first equation, we have \sol{$\frac{\ddot{x}}{\dot{x}}=\frac{2\dot{y}}{y}$}, integrating both sides with respect to $t$, we get $\ln|\dot{x}|=2\ln|y|+C_1\Rightarrow \sol{\dot{x}=cy^2}$.\par
    Since the geodesic has constant speed, we can assume unit speed without loss of generality. The length element squared is $ds^2=\frac{dx^2+dy^2}{y^2}=1$.Thus we have \sol{$\frac{\dot{x}^2+\dot{y}^2}{y^2}=1\Rightarrow \dot{x}^2+\dot{y}^2=y^2$}.\par
    Substitute $\dot{x}=cy^2$ into the speed equation, we have \sol{$(cy^2)^2+\dot{y}^2=y^2\Rightarrow c^2y^4+\dot{y}^2=y^2\Rightarrow \dot{y}^2=y^2(1-c^2y^2)$}.\par
    Thus $\frac{\d y}{\d x}=y^2(1-c^2y^2)\Rightarrow \pm\d x=\frac{cy}{\sqrt{1-c^2y^2}}\d y$.\par
    Take the integral, $\int\frac{cy}{\sqrt{1-c^2y^2}}\d y=\frac{-1}{2c}\int\frac{1}{\sqrt{1-c^2y^2}}\d(1-c^2y^2)=-\frac{1}{2c}\int\frac{1}{u^{-\frac{1}{2}}}\d u=-\frac{1}{c}\sqrt{1-c^2y^2}$.\par
    Therefore, we have \sol{$-\frac{1}{c}\sqrt{1-c^2y^2}=\pm(x-x_0)\Rightarrow \formu{(x - x_0)^2 + y^2 = \left(\frac{1}{c}\right)^2}$}
\end{enumerate}


\subsection{EXERCISE 5.}
\textit{Curve $c(t)=\left(t, \frac{1+t}{t}, \frac{1-t^{2}}{t}\right)$ for $t \in (0, \infty)$.}
\begin{enumerate}[label=(\roman*)]
    \item \textbf{Prove that $c$ is regular.}\par
    $\dot{c}(t)=(1, -t^{-2}, -t^{-2}-1)$.The first component is always 1, thus $\|\dot{c}(t)\|\neq 0$ for all $t$.
    \item \textbf{Compute the torsion of $c$.}\par
    $\ddot{c}(t)=(0, 2t^{-3}, 2t^{-3})$.Then \sol{$\dot{c} \times \ddot{c} = \begin{vmatrix} i & j & k \\ 1 & -t^{-2} & -1-t^{-2} \\ 0 & 2t^{-3} & 2t^{-3} \end{vmatrix}=2t^{-3}(1,-1,1)$}.\par
    Therefore, the curvature $\kappa$ is given by: \formula{$\kappa = \frac{\|\dot{c} \times \ddot{c}\|}{\|\dot{c}\|^3} = \formu{\frac{2\sqrt{3}/t^3}{\left( \frac{\sqrt{2(t^4+t^2+1)}}{t^2} \right)^3} = \frac{\sqrt{6} t^3}{(t^4+t^2+1)^{3/2}}}$}
    \item \textbf{Compute the torsion of $c$}.\par
    $\dddot{c} = (0, -6t^{-4}, -6t^{-4})$.$(\dot{c} \times \ddot{c}) \cdot \dddot{c}=2t^{-3}(1, -1, 1) \cdot -6t^{-4}(0, 1, 1) = -12t^{-7}(-1 + 1) = 0$.\par
    Thus, the torsion $\tau$ is \ex{$\tau = \frac{(\dot{c} \times \ddot{c}) \cdot \dddot{c}}{\|\dot{c} \times \ddot{c}\|^2} = \frac{0}{\|2t^{-3}(1, -1, 1)\|^2} = \formu{0}$}.
\end{enumerate}



\subsection{EXERCISE 6.}
\textit{Let $S$ be a compact surface with unit normal field $n: S \to S^2$.}

\begin{enumerate}[label=(\roman*)]
    \item \textbf{Prove that $n$ is surjective onto $S^2$.}\par
    Let $v \in S^2$ be an arbitrary unit vector in $\mathbb{R}^3$. We want to show that there exists a point $p \in S$ such that $n(p) = v$.

Consider the \textit{height function} $h_v: S \to \mathbb{R}$ defined by the dot product: $h_v(x) = x \cdot v, \quad \text{for all } x \in S$.\par
Since $S$ is compact and $h_v$ is continuous, $h_v$ must achieve its global maximum at some point $p \in S$.\par
Let $p$ be a point where $h_v$ attains its maximum. At a critical point of a function defined on a surface, the gradient of the restriction is zero. Geometrically, this means the tangent plane $T_p S$ is orthogonal to the gradient of the ambient linear function $f(x) = x \cdot v$. The gradient of $f$ is simply the vector $v$. Therefore, the vector $v$ is normal to the tangent plane $T_p S$.\par
This implies that the unit normal $n(p)$ is parallel to $v$. That is: $n(p) = v \quad \text{or} \quad n(p) = -v.$\par
To distinguish between these two, observe that $p$ is a global maximum. This means that for all $x \in S$, $x \cdot v \leq p \cdot v$. Geometrically, the entire surface $S$ lies in the half-space defined by the plane passing through $p$ with normal $v$. 

Assuming the standard convention that $n$ is the \textit{outward} pointing normal of a compact surface, the normal at the "highest" point in the direction $v$ must point in the direction of $v$ (out of the bounded region enclosed by $S$). Therefore, $n(p) = v$.

Since $v \in S^2$ was arbitrary, for every unit vector $v$, there exists a point $p \in S$ such that $n(p) = v$. Thus, $n$ is surjective.

    \item \textbf{Let $K^{+}(p) = \max\{0, K(p)\}$. Show that $\displaystyle \int_{S} K^{+} dA \geq 4\pi$.}\par
    Since $n$ is surjective, for every unit vector $v \in S^2$, there exists at least one point $p \in S$ such that $n(p) = v$. Each such point $p$ corresponds to a global maximum of the height function in the direction of $v$, and at these points, the Gaussian curvature $K(p)$ is non-negative. By considering the area of the sphere and the contributions from these points, we can show that the total positive curvature must be at least $4\pi$.
    \item \textbf{Determine the image of $n$ for the catenoid and the cylinder.}\par
    For the Catenoid, the normal vector is \ex{$n = (\mathrm{sech} v \cos u, \mathrm{sech} v \sin u, -\tanh v)$}.It covers the entire sphere except the north and south poles as $v$ varies over $(-\infty, \infty)$.\par
    For the Cylinder, the normal vector is \ex{$n = (\cos u, \sin u, 0)$}.It only covers the equator of the sphere as $u$ varies over $[0, 2\pi)$.
\end{enumerate}


\subsection{EXERCISE 7.}
\textit{Compact, oriented surface $S$ with mean curvature $H$ and Gauss curvature $K$.}
\begin{enumerate}[label=(\roman*)]
    \item \textbf{Prove that $H^{2} \geq K$.}\par
    By the definition, \sol{$H = \frac{k_1+k_2}{2}, K = k_1 k_2$}.$H^2-K=\left(\frac{\kappa_1+\kappa_2}{2}\right)^2-\kappa_1\kappa_2=\frac{(\kappa_1-\kappa_2)^2}{2}\geq 0$.
    \item \textbf{Show that $S$ cannot be a minimal surface.}\par
    If $S$ were a mininal surface, then $H=0$ for all points on $S$.From (i), we have $0=H^2\geq K\Rightarrow K\leq 0$ for all points on $S$.\par
    But $S$ is a compact surface, there exists at least one point $p$ such that $K(p)>0$.This contradicts the previous result that $K\leq 0$ everywhere on $S$.\par
    Therefore, $S$ cannot be a minimal surface.
    \item \textbf{Assume $S$ is not the sphere and that is does not have a smooth bijection to a sphere.Prove that then there are points on $S$ where the Gaussian curvature is positive, negative and zero.}\par
    Since $S$ is compact, there exists at least one point $p$ such that $K(p)>0$.\par
    By the Gauss-Bonnet theorem, \ex{$\int_S K dA = 2\pi \chi(S)$}.Since $S$ is not a sphere, its Euler characteristic $\chi(S) \leq 0$.Thus, the total Gaussian curvature $\int_S K dA \leq 0$.Thus, there must be points on $S$ where $K<0$ to ensure that the total curvature is non-positive.\par
    Since $S$ is connected and $K$ is a continuous function, by the Intermediate Value Theorem, there must be points on $S$ where $K=0$ as well.
\end{enumerate}

\section{Appendix}
\subsection{Integrals of Special Functions}
$+C$ and $\d x$ is omitted in the following formulas for brevity.
\begin{multicols}{2}
    \noindent\key{Trigonometric Functions}
    \begin{itemize}[leftmargin=*]
        \item \boldmath{$\int \tan x = -\ln|\cos x|$}
        \item \boldmath{$\int \cot x = \ln|\sin x|$}
        \item \boldmath{$\int \sec x = \ln|\sec x + \tan x|$}
        \item \boldmath{$\int \csc x = \ln|\csc x - \cot x|$}
    \end{itemize}

    \noindent\key{Inverse Trigonometric Functions}
    \begin{itemize}[leftmargin=*]
        \item \boldmath{$\int \arcsin x = x \arcsin x + \sqrt{1-x^2}$}
        \item \boldmath{$\int \arccos x = x \arccos x - \sqrt{1-x^2}$}
        \item \boldmath{$\int \arctan x = x \arctan x - \frac{1}{2}\ln(1+x^2)$}
        \item \boldmath{$\int \operatorname{arccot} x = x \operatorname{arccot} x + \frac{1}{2}\ln(1+x^2)$}
        \item \boldmath{$\int \operatorname{arcsec} x = x \operatorname{arcsec} x - \ln|x + \sqrt{x^2-1}|$}
        \item \boldmath{$\int \operatorname{arccsc} x = x \operatorname{arccsc} x + \ln|x + \sqrt{x^2-1}|$}
    \end{itemize}
    
    \noindent\key{Hyperbolic Functions}
    \begin{itemize}[leftmargin=*]
       \item \boldmath{$\int \sinh x = \cosh x$}
       \item \boldmath{$\int \cosh x = \sinh x$}
       \item \boldmath{$\int \tanh x = \ln (\cosh x)$}
       \item \boldmath{$\int \coth x = \ln |\sinh x|$}
       \item \boldmath{$\int \operatorname{sech} x = \arctan(\sinh x)$}
       \item \boldmath{$\int \operatorname{csch} x = \ln\left|\tanh \frac{x}{2}\right|$}
    \end{itemize}

    \noindent\key{Inverse Hyperbolic Functions}
    \begin{itemize}[leftmargin=*]
       \item \boldmath{$\int \operatorname{arsinh} x = x \operatorname{arsinh} x - \sqrt{x^2+1}$}
       \item \boldmath{$\int \operatorname{arcosh} x  = x \operatorname{arcosh} x - \sqrt{x^2-1}$}
       \item \boldmath{$\int \operatorname{artanh} x = x \operatorname{artanh} x + \frac{1}{2}\ln(1-x^2)$}
       \item \boldmath{$\int \operatorname{arcoth} x  = x \operatorname{arcoth} x + \frac{1}{2}\ln(x^2-1)$}
       \item \boldmath{$\int \operatorname{arsech} x = x \operatorname{arsech} x + \arcsin x$}
       \item \boldmath{$\int \operatorname{arcsch} x = x \operatorname{arcsch} x + \operatorname{arsinh} x$}
    \end{itemize}
    
    \noindent\key{Rational/Irrational Integrals}
    \begin{itemize}[leftmargin=*]
        \item \boldmath{$\int \frac{1}{\sqrt{x^2+a^2}} = \ln |x+\sqrt{x^2+a^2}|$}
        \item \boldmath{$\int \frac{1}{\sqrt{x^2-a^2}} = \ln |x+\sqrt{x^2-a^2}|$}
        \item \boldmath{$\int \frac{1}{\sqrt{a^2-x^2}} = \arcsin \frac{x}{|a|}$}
        \item \boldmath{$\int \frac{1}{x^2+a^2} = \frac{1}{|a|} \arctan \frac{x}{|a|}$}
        \item \boldmath{$\int \frac{1}{x^2-a^2} = \frac{1}{2|a|} \ln \left|\frac{x-a}{x+a}\right|$}
        \item \boldmath{$\int \sqrt{x^2+a^2} = \frac{x}{2}\sqrt{x^2+a^2} + \frac{a^2}{2}\ln|x+\sqrt{x^2+a^2}|$}
        \item \boldmath{$\int \sqrt{x^2-a^2} = \frac{x}{2}\sqrt{x^2-a^2} - \frac{a^2}{2}\ln|x+\sqrt{x^2-a^2}|$}
        \item \boldmath{$\int \sqrt{a^2-x^2} = \frac{x}{2}\sqrt{a^2-x^2} + \frac{a^2}{2}\arcsin \frac{x}{|a|}$}
        \item \boldmath{$\int \frac{1}{\sqrt{(x^2+a^2)^3}} = \frac{x}{a^2\sqrt{x^2+a^2}}$}
        \item \boldmath{$\int \frac{1}{\sqrt{(x^2-a^2)^3}} = -\frac{x}{a^2\sqrt{x^2-a^2}}$}
    \end{itemize}
\end{multicols}

\subsection{Hyperbolic Trigonometric Identities}
    \begin{multicols}{2}
        \begin{itemize}[leftmargin=*]
    \item \boldmath{$\sinh x=\frac{e^x-e^{-x}}{2},\cosh x=\frac{e^x+e^{-x}}{2}$}
    \item \boldmath{$\tanh x=\frac{\sinh x}{\cosh x}=\frac{e^x-e^{-x}}{e^x+e^{-x}}$}
    \item \boldmath{$\cosh^2 x - \sinh^2 x = 1$}
    \item \boldmath{$\operatorname{sech}^2 x + \tanh^2 x = 1$}
    \item \boldmath{$\coth^2 x - \operatorname{csch}^2 x = 1$}
    \item \boldmath{$\sinh (x \pm y) = \sinh x \cosh y \pm \cosh x \sinh y$}
    \item \boldmath{$\cosh (x \pm y) = \cosh x \cosh y \pm \sinh x \sinh y$}
    \item \boldmath{$\tanh (x \pm y) = \frac{\tanh x \pm \tanh y}{1 \pm \tanh x \tanh y}$}
    \item \boldmath$\operatorname{arcsinh} x=\ln{x+\sqrt{x^2+1}}$
    \item \boldmath$\operatorname{arcosh} x=\ln{x+\sqrt{x^2-1}}$
    \item \boldmath$\operatorname{artanh} x=\frac{1}{2}\ln{\frac{1+x}{1-x}}$
\end{itemize}
    \end{multicols}

\subsection{Translations of common Curves and Surfaces}

\begin{multicols}{3}
    \begin{itemize}[leftmargin=*]
        \item Cylinder: 圆柱面
        \item Torus： 环面
        \item Catenoid: 双曲面(悬链面)
        \item Helicoid: 螺旋面
        \item Ellipsoid: 椭球面
        \item elliptic paraboloid: 椭圆抛物面
        \item hyperbolic paraboloid: 双曲抛物面
        \item hyperboloid of one sheet: 单叶双曲面
        \item hyperboloid of two sheets: 双叶双曲面
        \item Saddle surface: 马鞍面
        \item Cone: 锥面
        \item Helix： 螺旋线
        \item Cycloid: 摆线
        \item Ellipse: 椭圆
        \item Parabola: 抛物线
        \item Hyperbola: 双曲线
    \end{itemize}
\end{multicols}


\end{document}
